\section{Typed definitions and identity criteria}
\label{app:typed}

This appendix expands Definition~\ref{def:package} sufficiently to distinguish
the relations used in the paper.  It is a metalanguage for proof interfaces,
not an axiomatization of all physical theories.

\subsection{Components}

The carrier $\mathcal K$ is a measurable, topological, smooth, algebraic, or
Hilbert-space object as required by the application.  The redundancy
groupoid $\mathcal G\rightrightarrows\mathcal K$ specifies which carrier
presentations are identified for the physical claim.  An invariant algebra
may be defined directly, or by descent:
\[
 \Aphys=\{A\in\mathcal A_{\mathrm{cov}}:
  A(g\cdot k)=A(k)\ \text{for all admissible }g,k\}.
\]
The topology on $\Aphys$ is part of the package because density and
continuity depend on it.

A state is a normalized positive functional
$\omega:\Aphys\to\mathbb C$.  If $\Aphys$ is not complete, the package states
whether positivity and continuity refer to the algebraic, $C^*$, nuclear, or
distributional topology.  The dynamical component is one of the following
typed forms:
\begin{enumerate}[label=(\alph*),itemsep=2pt]
 \item a one-parameter automorphism or endomorphism group $\tau_t$;
 \item a positive Euclidean semigroup $\e^{-tH}$ after reconstruction;
 \item a family of $n$-point correlation functionals;
 \item a Markov kernel used only as an auxiliary sampling dynamics; or
 \item an RG map on measures/actions and retained states.
\end{enumerate}
The type is recorded precisely because (d), for example, has no automatic
spectral interpretation as (b).

The limit record $\mathcal L$ includes a directed set $I$, maps from the index
to cutoff and physical volume, observable-renormalization maps, admissible
cofinal nets, and a target topology.  A property predicate $Q\in\mathcal Q$
is a tuple
\[
 Q=(\text{bearer},\text{formula},\text{domain},\text{units},
     \text{quantifier prefix}).
\]
For example, a finite-lattice transfer eigenvalue and a continuum Hamiltonian
gap are different predicates even when both are represented by positive
numbers.

\subsection{Morphism data}

A morphism $F:\Rep_0\to\Rep_1$ includes:
\begin{enumerate}[label=(\roman*),itemsep=2pt]
 \item a domain $\mathcal D_1\subseteq\Aphys^{(1)}$ and observable pullback
 $F^*: \mathcal D_1\to\Aphys^{(0)}$;
 \item a compatible state map on a declared subset, when such a map exists;
 \item support, scale, and unit transformations;
 \item an expectation-fidelity modulus as in Equation~\eqref{eq:fidelity};
 \item exact, forgotten, recovered, and uniform certificate entries; and
 \item a statement of whether $F$ is an equivalence, quotient, embedding,
 regularization, coarse graining, continuum comparison, reconstruction, or
 heuristic correspondence.
\end{enumerate}

The identity morphism has identity pullbacks and pushforwards, zero error,
every declared property exact, no forgotten channel, and the full index
domain.  An exact equivalence requires inverse morphisms whose composites are
identities on the stated physical algebras and states, plus compatibility of
interpretation.  A quotient is exact only for observables descending through
the equivalence relation.  A coarse graining is not assigned an inverse unless
a separate recovery map is proved on a specified subspace.

\subsection{Recovery modes}

Four recovery modes are kept distinct.

\begin{description}[style=nextline,leftmargin=8mm]
 \item[Exact section.] A map $S$ with $F\circ S=\id$ or $S\circ F=\id$ on
 the stated domain.
 \item[Observable reconstruction.] A map from target data back to source
 expectation values on a chosen core, with an error modulus.
 \item[Asymptotic recovery.] A net of reconstructions whose error converges
 to zero in the declared topology and quantifier order.
 \item[External theorem.] A theorem, such as OS reconstruction, that creates
 a target bearer from data satisfying hypotheses not manifest as an inverse
 map.
\end{description}

The recovery mode is part of the certificate.  Stating that information is
``implicit'' without one of these mechanisms does not close a proof
obligation.

\subsection{Uniformity domains}

A uniformity record is a formula, not a tag.  A typical entry has the form
\[
 \exists m>0\ \exists a_0>0\ \forall a<a_0\
 \forall L\geq L_0(a)\ \forall A,B\in\mathcal D_{a,L}\
 \forall t\geq0:\quad \Phi(a,L,A,B,t;m).
\]
Changing the order can change the claim.  The record also states whether
$L_0$ depends on $a$, whether supports remain at fixed physical size, and
whether the bound is pointwise or distributional in $t$.
